\begin{proposition}\hypertarget{tag-01G0}{}\label{cdp:st:D.1.2}
\StatementRecord{cdp:st:D.1.2}{proved}{01G0}{\SeriesRef{cdp:eq:D1}}

Let $f:Y\to C$ be a proper holomorphic surjection from a connected
compact complex threefold to a compact Riemann surface, with
connected fibres. Let $\Delta$ be a nonempty finite set containing
the critical values, $s=|\Delta|$, $g=g(C)$,
$D=f^{-1}(\Delta)_{\rm red}$, and let $r$ be the number of
irreducible components of $D$. Write $F$ for a smooth fibre and
\[
t'=\dim_{\Q}H_1(F;\Q)_{\pi_1(C\setminus\Delta)}
   =\dim_{\Q}H^1(F;\Q)^{\pi_1(C\setminus\Delta)}.
\]
For $b_2(Y)=0$ the exact component count is
\[
\boxed{\quad r=2g+s-1+t'-b_1(Y).\quad}
\]
Under CDP20's integral assumptions
$H^1(Y;\Z)=H^2(Y;\Z)=0$, the base is $\PP^1$ and
\[
r=s-1+t',\qquad
H_1(F;\Z)_{\pi_1(C\setminus\Delta)}\text{ is free abelian}.
\]
The last assertion concerns the indicated quotient of $H_1(F;\Z)$;
it does not assert that $H_1(F;\Z)$ itself is torsion-free.
\end{proposition}



The rational component-count formula is the correction in \citet[Section~10.5, p.~91]{AC2026}. The integral torsion-free conclusion under the displayed integral hypotheses is an additional clause; its full argument is included here.

\begin{proof}

\begin{enumerate}[label={$\langle1\rangle$\arabic*.},leftmargin=*]
\item Put $U=Y\setminus D$ and $C^\circ=C\setminus\Delta$.
The map $U\to C^\circ$ is a differentiable fibre bundle. Since
$\pi_2(C^\circ)=0$, its homotopy sequence is
\[
1\longrightarrow\pi_1(F)\longrightarrow\pi_1(U)
 \longrightarrow\pi_1(C^\circ)\longrightarrow1.
\]
The last group is free of rank $2g+s-1$. Lifting its free generators
gives a section of this extension. Proposition \CompendiumRef{II}{cdp:st:D.1.1} therefore gives
\[
H_1(U;\Z)\simeq
\Z^{2g+s-1}\oplus H_1(F;\Z)_{\pi_1(C^\circ)}.       \tag{D1}\hypertarget{tag-019T}{}\label{cdp:eq:D1}
\]
The second equality defining $t'$ follows by duality: the dual of
the coinvariant quotient is precisely the space of invariant
linear functionals on $H_1(F;\Q)$.

\item The cohomology sequence with compact supports for the open inclusion
$U\subset Y$ contains
\[
H^4(Y;\Q)\longrightarrow H^4(D;\Q)
\longrightarrow H_c^5(U;\Q)
\longrightarrow H^5(Y;\Q)\longrightarrow H^5(D;\Q)=0.
\]
The top-degree cohomology of the compact pure complex surface $D$
is $\Q^r$. One can see this by deleting the singular and intersection
loci, whose real dimensions are at most two, and taking compactly
supported top cohomology of its $r$ oriented smooth components.
Poincar\'e duality identifies $H_c^5(U;\Q)$ with $H_1(U;\Q)$ and
$H^5(Y;\Q)$ with $H_1(Y;\Q)$. Since $H^4(Y;\Q)=0$, taking
dimensions and using \eqref{cdp:eq:D1} proves the formula.

\item For the integral refinement, the universal coefficient theorem shows
first that $H_1(Y;\Z)$ is finite, using $H^1(Y;\Z)=0$, and then that
it is zero, using
$\operatorname{Ext}^1(H_1(Y;\Z),\Z)\subset H^2(Y;\Z)=0$.
The group $H^4(Y;\Z)$ can be finite and nonzero. Its homomorphism
to $H^4(D;\Z)=\Z^r$ is nevertheless zero. The same sequence now
gives $H_1(U;\Z)\simeq\Z^r$, since $H^5(Y;\Z)=H_1(Y;\Z)=0$.
Equation \eqref{cdp:eq:D1} then proves that the coinvariant summand is free.
Finally $H^1(C;\Q)\hookrightarrow H^1(Y;\Q)$ by the Leray edge
map, so $g=0$.

\end{enumerate}
\end{proof}